\section{Named model backends: a scientific contract above the generation interface}
\label{sec:exec24-backend-census}

A bank considering an assistant for causal study design must choose both the
model service and the execution environment through which it is used. A useful
comparison asks what each environment preserves when a question becomes a
diagram, a tool request and a reviewed answer. It also asks what happens when
the output is partial, a tool fails, the request is retried or the budget is
exhausted. An API-shaped response is only one intermediate artifact. The
scientific object remains the question, information clock, population,
intervention, assumptions, graph and justified inference status.
\citep{exec24_backend_reference}

The selected primary profiles are the official Anthropic Python SDK 1.11.0,
Google GenAI Python SDK 2.28.0, the Ollama Python client, vLLM and llama.cpp.
Each repository revision is pinned no later than 7 October 2026. Ollama's
package version is derived from version-control metadata; the inspected
vLLM tree likewise does not establish a release wheel version. The llama.cpp
build declares base version 0.6.0 with development mode enabled by default.
These distinctions prevent a commit, a package release, a model name and a
deployed service from being treated as the same identity.
\citep{exec24_backend_anthropic,exec24_backend_google,exec24_backend_ollama,
exec24_backend_vllm,exec24_backend_llamacpp}

The comparison is deliberately bounded. Official OpenAI guide originals
could not be acquired through the permitted routes; repeated HTTP 403
responses remain an omission, not evidence that the service lacks a feature.
Hosted-service entitlements, immutable service revisions, complete model
catalogues, regional behavior and actual prices are not established by the
five inspected source profiles. Azure and Bedrock deployments, additional
model hosts and other local engines require their own applicable evidence.
The following choices therefore supply a prototype route while leaving the
broader backend census open.\citep{exec24_backend_reference}

\subsection{Separate syntax, object validation and scientific validity}

Let $x$ denote a generated representation and let $S(x)$ mean that it satisfies
a declared serialization schema. Let $G(x)$ mean that it denotes an admissible
graph and $I(x)$ mean that the intervention query is justified under the
declared causal assumptions. These are different predicates. A conventional
edge-list schema can admit the three edges $A\to B$, $B\to C$ and $C\to A$:
every endpoint is a valid named variable and every field has the expected
type, but the graph is cyclic. An independent topological check rejects it
when the contract requires a DAG. Even an acyclic graph does not establish
that it represents the bank's assignment process. The product must therefore
preserve separate schema, graph and identification decisions rather than
promoting $S(x)$ directly to $I(x)$.
\citep{exec24_backend_reference}

The Anthropic SDK illustrates another level of this distinction. Its selected
output configuration contains a JSON-schema format field. Separately, the
client's internal \texttt{\_strict\_response\_validation} option defaults to
false. The response-processing routine invokes a type validator when that
option is enabled and otherwise constructs the target object; a reported
validation error is wrapped as an API response-validation error. These are
client object-handling branches, not a causal validity oracle or proof about
every service-side generation path. A proposed adapter keeps the requested
schema, returned content and observed validation status separately.
\citep{exec24_backend_anthropic}

Google's selected response conversion is more explicit about the requested
Python representation. A Pydantic model invokes its JSON validator; caught
JSON or validation errors leave the parsed result unset. A dictionary or
SDK Schema branch instead applies JSON parsing to the returned text, without
an additional local full-schema validator in that branch. A caller must
not interpret a returned response object, or merely successful JSON parsing,
as proof that all requested schema constraints hold. These are inspected
client branches, not measurements of service violation rates.
\citep{exec24_backend_google}

For vLLM, the pinned documentation distinguishes schema-derived decoding
from structural-tag parsing. In the latter route, constraining the tool-call
envelope and enforcing an individual tool's argument schema are separate.
The documented per-tool strict flag, server strictness setting and global
enforcement switch affect that route. General named-function prose is
qualified by this later parser-specific explanation; it is unsafe to repeat
the broad statement as an unconditional guarantee across all parsers.
The structured-output configuration also chooses a backend automatically
by default and describes that choice as release-dependent. A prototype must
pin the parser and decoding backend as well as the model and server revision.
\citep{exec24_backend_vllm}

llama.cpp exposes grammar and JSON-schema generation controls and directs
readers to supported-feature tests. That is a useful constrained-generation
boundary, but neither the parameter name nor a syntactically valid output
establishes full JSON Schema coverage, a valid causal graph or identification.
The proposed common validator runs after every backend and retains the
backend's original output for diagnosis. Unsupported constraints must be
reported, not silently weakened to make a request run.
\citep{exec24_backend_llamacpp,exec24_backend_reference}

\subsection{Tool requests, tool execution and scientific responsibility}

The Google SDK documents automatic function calling when a Python function
is supplied directly. The inspected loop counts remote model calls, retrieves
function-call parts and invokes the registered functions through a helper.
At the last permitted remote call it returns the model's function calls
without executing them, because no remote call remains for their results.
The helper can visit several function-call parts in a single response and
can convert a local exception into an error-bearing function response.
Consequently a remote-call limit is not a limit on the number, duration or
economic consequences of local function invocations.
\citep{exec24_backend_google}

In an authored example, two model responses are permitted. The first contains
three executable tool calls; the final response contains two more that remain
unexecuted under the stipulated last-call rule. There are two remote calls,
three executed local functions and two unresolved requests. With hypothetical
model costs of 0.20 and 0.30 USD and tool costs of 2, 5 and 3 USD, expenditure
is 10.50 USD. Reporting only two calls loses both execution and cost meaning.
For a credit-policy assistant, the adapter should reserve separate budgets
for model attempts, scientific tools, human review, wall time and monetary
expenditure before allowing an action.
\citep{exec24_backend_reference}

The selected Ollama \texttt{chat} path accepts tools, converts their
declarations, sends the request and returns a typed response or a stream of
typed chunks. That path does not itself implement the Google-style local
function-execution loop. Applications using the two interfaces therefore
need different orchestration policies even when both accept a tools argument.
For the proposed library, generation may suggest a causal operation, while
the controller checks its scientific context and permission before execution.
A model request for estimation cannot create identification that the causal
engine has refused.\citep{exec24_backend_ollama,exec24_backend_reference}

\subsection{Completion, refusal and observation deadlines}

The Anthropic message type distinguishes a natural end, token limit,
specified stop sequence, tool request, paused turn, refusal and context-window
limit. Its documentation says that a streaming message begins without a final
stop reason. Google's finish-reason type likewise separates ordinary stopping,
token limits, safety filtering, malformed function calls and other outcomes.
Ollama's response model carries \texttt{done}, \texttt{done\_reason}, token
counts and durations expressed in nanoseconds. A common adapter should
retain native states before mapping them to its own completion, partial,
refused or failed categories. Provider-policy refusal is distinct from an
engine's scientific refusal to identify an effect.
\citep{exec24_backend_anthropic,exec24_backend_google,exec24_backend_ollama}

Time budgets need equally explicit units and scope. The Anthropic constants
set a ten-minute default HTTP timeout profile and two retries; the inspected
retry routine considers selected HTTP errors and server retry hints. Google's
HTTP options describe timeout in milliseconds and maximum attempts including
the initial request. Ollama's client explicitly defaults timeout to none.
These observations do not establish a whole-business-workflow deadline.
Retries, model generation, local tools, queueing and result review require a
separate end-to-end clock and a declared policy for unresolved remote work.
\citep{exec24_backend_anthropic,exec24_backend_google,exec24_backend_ollama}

The distinction has exact consequences even without calling a provider. Two
retries permit three request attempts; three stipulated four-second attempts
and two twelve-second retry waits total 36 seconds. Calling this a
four-second operation because each observed request took four seconds would
omit most elapsed time. This example is accounting under a specified sequence,
not a measured latency or a universal upper bound.
\citep{exec24_backend_reference}

The selected llama.cpp source makes the prediction-time qualification
concrete. The inspected stopping branch tests elapsed generation time inside
the branch that sends a newline. Its documented limit concerns prediction,
not queueing and prompt processing. With hypothetical queueing of five
seconds, prefill of three seconds and the next relevant newline at eight
seconds of generation, a two-second generation setting can be observed at
sixteen seconds end to end, assuming no earlier stop. This is a conditional
control-flow witness, not a benchmark. The inspected Cancel task releases a
matching slot; that operation alone does not reverse a previously executed
external tool or establish every transport-disconnection behavior.
\citep{exec24_backend_llamacpp,exec24_backend_reference}

\subsection{Compatible endpoints and local deployment are separate choices}

vLLM documents several OpenAI-compatible routes while identifying parameter
exceptions and model-dependent tool behavior. Its server can apply sampling
defaults from a model repository's generation configuration unless configured
otherwise. llama.cpp exposes compatible routes and warns that prompt-cache
reuse can yield nonidentical logits across batch sizes on some backends.
Thus matching an endpoint and seed does not establish identical generation
semantics or bitwise reproducibility. A reproducibility record binds model
weights, tokenizer, chat template, quantization, decoding backend, sampling
configuration, cache and hardware conditions.
\citep{exec24_backend_vllm,exec24_backend_llamacpp}

The direct source notices differ from the rights of those artifacts.
Anthropic's Python SDK, the Ollama client and llama.cpp carry MIT notices;
Google GenAI and vLLM declare Apache 2.0. These observations concern selected
source distributions. They do not provide hosted-service access, model-weight
rights, training-data permissions, a redistributable dependency environment
or a bank support agreement. The proposed first comparison references these
implementations, conditionally integrates an original client or local worker
after qualification, and rejects an unmotivated wholesale port. Rewriting a
client would preserve the missing service and scientific evidence while adding
a second maintenance burden.
\citep{exec24_backend_anthropic,exec24_backend_google,exec24_backend_ollama,
exec24_backend_vllm,exec24_backend_llamacpp,exec24_backend_reference}

\subsection{A public comparison that can justify a prototype}

The first experiment should use public or generated business-to-causal
questions with explicit target, graph and refusal oracles. The same logical
cases, evidence and rubrics should be paired across backends; repeated attempts
measure within-case variation rather than creating new independent cases.
Record all submissions, completed responses, parseable objects, scientifically
acceptable answers, appropriate refusals and unresolved failures.
\citep{exec24_backend_reference}

For illustration, backend A correctly answers 45 of 50 completed cases from
100 submissions, while B correctly answers 72 of 90 completed cases from
the same number submitted. Conditional agreement favors A, 90 versus
80 percent; correct answers per submitted case favor B, 72 versus 45 percent.
At hypothetical total costs of 150 and 180 USD, cost per correct answer is
$10/3$ versus $5/2$ USD. These invented data show why the evaluation needs
both conditional quality and coverage; they rank no real provider.
\citep{exec24_backend_reference}

The scientific owner specifies the cases, loss and acceptance standard;
the adapter owner preserves state, units and artifacts; the platform owner
qualifies resource and rights conditions; and the evaluation owner accounts
for attempts, exclusions and full cost. Broader hosted evidence and native
trials remain necessary before a named backend is selected for institutional
use. The comparison makes the next public experiment concrete without
claiming that client code, constrained generation or a completed request
constitutes a qualified causal product.
\citep{exec24_backend_reference}
